\texttt{ET-01} & ET & Qual a dose de isoniazida no esquema 3HP para adultos? & No esquema 3HP, adultos (>14 anos, ≥30kg) tomam 900mg de isoniazida por semana, uma vez por semana, durante 3 meses (12 doses). O esquema prevê ajuste de dose por peso corporal. \\ \hline
\texttt{ET-02} & ET & Qual a duração do esquema 3HP e quantas doses são necessárias? & O tratamento 3HP estará completo quando ocorrer a tomada de 12 doses de isoniazida + rifapentina por 12 semanas. Dependendo do caso, esse prazo pode ser prorrogado para 15 semanas. \\ \hline
\texttt{ET-03} & ET & Qual a dose de rifampicina no esquema 4R para adultos e por quanto tempo? & Rifampicina 600 mg/dia, indicada para adultos ≥ 50 kg, administrada diariamente por 4 meses, totalizando 120 doses; que podem ser completadas no período de 4 a 6 meses, conforme adesão ao tratamento. \\ \hline
\texttt{ET-04} & ET & Qual a dose de isoniazida no esquema 6H para crianças? & Esquemas pediátricos incluem isoniazida (10 mg/kg/dia; dose máxima: 300 mg/dia) administrada diariamente por 6 meses (totalizando 180 doses, a serem completadas em 6 a 9 meses). \\ \hline
\texttt{ET-05} & ET & Qual esquema é preferencial para hepatopatas e idosos acima de 50 anos? & A rifampicina é o tratamento de escolha para hepatopatas e pessoas acima de 50 anos. \\ \hline
\texttt{ET-06} & ET & Qual a dose pediátrica de isoniazida no esquema 3HP para crianças de 10 a 15 kg? & No esquema 3HP pediátrico (crianças de 2 a 14 anos), a dose de isoniazida para crianças de 10 a 15 kg é de 300mg por semana. Para rifapentina, a dose é também 300mg por semana para esse peso. \\ \hline
\texttt{ET-07} & ET & Qual a dose de isoniazida para adultos no esquema 6H? & Para adultos, a dose de isoniazida nos esquemas 6H e 9H é de 5 a 10 mg/kg/dia, com dose máxima de 300 mg/dia. O esquema 6H contempla 180 doses (6 a 9 meses) e o 9H contempla 270 doses (9 a 12 meses). O esquema 9H protege mais que o 6H e deve ser preferido sempre que houver boa adesão. \\ \hline
\texttt{IT-01} & IND & Quem deve receber tratamento da ILTB independente do resultado da PT ou IGRA? & Alguns grupos devem ser tratados sem PT e sem IGRA, como contatos de alto risco com TB bacilífera. O Quadro 1 do Protocolo de Vigilância da ILTB especifica as indicações de tratamento, incluindo grupos que devem tratar sem PT/IGRA, grupos que tratam com PT ≥ 5mm ou IGRA positivo, e grupos que tratam com PT ≥ 10mm ou IGRA positivo. \\ \hline
\texttt{IT-02} & IND & Quando é indicado tratar ILTB em pessoa com alteração radiológica fibrótica pulmonar? & Todos os indivíduos com PPD maior ou igual a 5mm e IGRA positivo e que tenham alterações radiológicas pulmonares fibróticas sugestivas de sequelas de TB, independe de ter alguma doença ou estar em tratamento, também têm indicação de tratamento da TBi. \\ \hline
\texttt{IT-03} & IND & Qual a indicação de tratamento da ILTB em paciente em pré-transplante? & Estão incluídos na indicação de tratamento da TBi os pacientes com PPD maior ou igual a 5mm ou IGRA positivo em pré-transplante em uso de terapia imunossupressora. \\ \hline
\texttt{IT-04} & IND & Quando tratar ILTB em paciente com prednisona? & Devem receber tratamento para TBi indivíduos com PPD maior ou igual a 5mm e IGRA positivo, que estejam em uso ou iniciarão terapia com dose acima de 15 mg/dia de prednisona por período superior a um mês. \\ \hline
\texttt{IT-05} & IND & Quando é indicado tratar ILTB em pessoas em contato com TB? & Recomenda-se o tratamento da infecção latente pelo M. tuberculosis (ILTB) para pessoas infectadas, identificadas por meio da prova tuberculínica (PT) ou Interferon-Gamma Release Assays (IGRA), quando apresentam risco de desenvolver TB. Antes do tratamento deve-se afastar definitivamente a tuberculose ativa. \\ \hline
\texttt{PE-01} & PE & Posso tratar ILTB em gestante HIV negativa? & Em gestantes HIV negativas, a recomendação é postergar o tratamento da ILTB para após o parto. \\ \hline
\texttt{PE-02} & PE & Quando tratar ILTB em gestante HIV positiva? & Em gestantes HIV positivas, preconiza-se o tratamento após o 3º mês de gestação. \\ \hline
\texttt{PE-03} & PE & Como manejar paciente em uso de anti-TNF com diagnóstico de ILTB? & Indivíduos com PPD ≥ 5mm e IGRA positivo que estejam em uso ou iniciarão terapia biológica com ação anti-TNF alfa têm indicação de tratamento para TBi. O tratamento deve ser iniciado 30 dias antes do uso da terapia imunossupressora. Após o diagnóstico durante o tratamento da DII, recomenda-se tratamento da TBi além de vigilância e monitoramento. \\ \hline
\texttt{PE-04} & PE & Qual esquema de tratamento da ILTB deve ser evitado em gestantes? & O esquema 3HP (Rifapentina + Isoniazida) não está indicado para gestantes. As indicações do 3HP excluem gestantes entre as populações para as quais o esquema pode ser utilizado. \\ \hline
\texttt{PE-05} & PE & Qual a dose de rifampicina para crianças no esquema 4R? & No esquema pediátrico 4R, a rifampicina é administrada na dose de 15 mg/kg/dia, com dose máxima de 600 mg/dia, administrada diariamente por 4 meses. \\ \hline
\texttt{PE-06} & PE & Qual o risco aumentado de TB doença em pacientes com DII em uso de anti-TNF? & Estudo realizado em nosso meio demonstrou aumento significativo do risco de TB doença em pacientes com DII em tratamento. O uso de tiopurinas aumentou o risco em 5,85 vezes, anti-TNF em 3,93 vezes e a combinação azatioprina + anti-TNF em 9,03 vezes. \\ \hline
\texttt{PE-07} & PE & Qual esquema usar para paciente com HIV? & Para PVHIV, o esquema preferencial é a isoniazida (9H), pois o esquema 4R (rifampicina) está contraindicado em PVHIV em uso de inibidores de protease e integrase. O esquema 9H (270 doses, 9 a 12 meses) protege mais que o 6H. PVHIV com contagem de CD4 ≤ 350 células/µL têm indicação de tratamento para ILTB independentemente do resultado da PT ou IGRA. \\ \hline
\texttt{DI-01} & DI & Qual o ponto de corte da prova tuberculínica para indicar tratamento da ILTB em pacientes em uso de anti-TNF? & Devem receber tratamento para TBi indivíduos com PPD maior ou igual a 5 mm e IGRA positivo, que estejam em uso ou iniciarão terapia biológica com ação anti-TNF alfa ou uso de dose acima de 15 mg/dia de prednisona por período superior a um mês. \\ \hline
\texttt{DI-02} & DI & Qual exame é superior para diagnóstico da ILTB: IGRA ou PPD? & Não há superioridade do IGRA sobre o PPD e vice-versa. O PPD sofre influência da vacina BCG até dois anos depois de sua aplicação. No caso do IGRA, existe a possibilidade de resultado indeterminado, especialmente em condições de moradia em região endêmica, tratamento com corticoide ou imunossupressor, atividade de doença, idade avançada, baixos níveis de albumina, linfopenia e altos níveis de proteína C reativa. \\ \hline
\texttt{DI-03} & DI & O que deve ser feito antes de iniciar o tratamento da ILTB? & Antes de efetuar o tratamento da ILTB deve-se: (a) afastar definitivamente a tuberculose ativa e (b) realizar a notificação por meio da Ficha de notificação das pessoas em tratamento da ILTB. O diagnóstico de ILTB deve ser realizado por meio dos testes da PT ou do IGRA. \\ \hline
\texttt{DI-04} & DI & Quando o resultado do IGRA for indeterminado, o que fazer? & Quando o resultado do IGRA for indeterminado, o exame deve ser repetido o mais breve possível. Se o resultado for mantido, é recomendado considerar o tratamento de TBi. \\ \hline
\texttt{DI-05} & DI & Qual a conduta para excluir TB ativa antes de tratar ILTB? & Para excluir TB ativa antes de tratar ILTB, deve-se realizar: anamnese, exame clínico, laboratorial e radiografia de tórax (por vezes tomografia computadorizada de tórax). É importante excluir todas as formas clínicas da doença (pulmonar/laríngea e extrapulmonar), além de considerar vínculos epidemiológicos e exames bacteriológicos. \\ \hline
\texttt{MO-01} & MO & Com que frequência devem ser realizadas as avaliações clínicas no esquema 3HP? & No esquema com isoniazida e rifapentina (3HP) estão indicadas avaliações mensais para avaliação clínica, seguimento da adesão ao tratamento e de eventos adversos. \\ \hline
\texttt{MO-02} & MO & Quando devo suspender o tratamento da ILTB por hepatotoxicidade? & Durante o tratamento da ILTB, a hepatite aguda medicamentosa (CID K71) é uma das causas de interesse para investigação e registro. O tratamento deve ser suspenso diante de sinais de hepatotoxicidade clinicamente significativa, como elevação de transaminases com sintomas. \\ \hline
\texttt{MO-03} & MO & O que considerar durante o acompanhamento do tratamento da ILTB além do tempo? & Durante o tratamento de TBi é importante levar em consideração o número de doses utilizadas e não apenas o tempo de tratamento. \\ \hline
\texttt{MO-04} & MO & Quais informações devem ser registradas ao final de cada consulta de acompanhamento? & Ao final da consulta de seguimento, deve-se: agendar consulta de seguimento para monitorar o progresso do tratamento, providenciar uma forma de o paciente entrar em contato com um profissional de saúde caso surjam problemas, organizar consultas em grupo se houver, registrar o que aconteceu durante a consulta, e encaminhar para serviços sociais existentes. \\ \hline
\texttt{MO-05} & MO & Com que frequência pacientes com DII em uso de imunossupressor devem ser rastreados para ILTB? & Pacientes com DII em uso contínuo de terapia com ação imunossupressora devem realizar rastreio anual para TBi por meio do PPD e/ou IGRA. Além disso, o consenso recomendou que pacientes sem tratamento prévio de TB realizem exames para triagem de TBi anuais durante o período de 3 anos e, em caso de mudança da terapia, a triagem está recomendada por um período de mais 3 anos. \\ \hline
\texttt{IM-01} & IT & Rifampicina tem interação com contraceptivos orais? & A rifampicina interage com contraceptivos orais, reduzindo sua eficácia por indução enzimática. É recomendado utilizar método alternativo de contracepção durante o uso da rifampicina. \\ \hline
\texttt{IM-02} & IT & Qual a principal preocupação ao usar rifampicina em pessoas vivendo com HIV em uso de antirretrovirais? & A rifampicina é um potente indutor enzimático que pode reduzir significativamente os níveis plasmáticos de antirretrovirais, especialmente inibidores de protease e inibidores não nucleosídeos da transcriptase reversa. Deve-se verificar interações medicamentosas antes de iniciar rifampicina em PVHIV. \\ \hline
\texttt{IM-03} & IT & Isoniazida tem interação com fenitoína? & A isoniazida interage com fenilhidantoína (fenitoína), causando maior hepatotoxicidade. A recomendação do Ministério da Saúde é evitar o uso concomitante. Quando ambos os fármacos forem necessários, deve-se monitorar sintomas e enzimas hepáticas conforme indicado. \\ \hline
\texttt{IM-04} & IT & Quanto tempo antes do início de terapia imunossupressora deve ser iniciado o tratamento da ILTB? & O Ministério da Saúde recomenda que o tratamento da TBi seja iniciado 30 dias antes do uso da terapia imunossupressora. \\ \hline
\texttt{IM-05} & IT & Qual o prazo recomendado para aguardar antes de iniciar tratamento imunossupressor após diagnóstico de ILTB, segundo o consenso europeu? & O Consenso Europeu que aborda infecções em DII recomenda aguardar 4 semanas para início do tratamento imunossupressor após o diagnóstico de ILTB, exceto em casos de maior urgência clínica após aconselhamento especializado. \\ \hline
\texttt{EA-01} & EA & Quais são os principais efeitos adversos da isoniazida? & Os principais efeitos adversos da isoniazida incluem: hepatotoxicidade (hepatite aguda medicamentosa), neuropatia periférica (pode ser prevenida com piridoxina), reações de hipersensibilidade. A hepatotoxicidade é a reação adversa mais grave e requer monitoramento das transaminases. \\ \hline
\texttt{EA-02} & EA & Como prevenir a neuropatia periférica causada pela isoniazida? & A neuropatia periférica causada pela isoniazida pode ser prevenida com a suplementação de piridoxina (vitamina B6). Em gestantes HIV negativas com indicação de tratamento, a isoniazida deve ser acompanhada de suplementação de piridoxina. \\ \hline
\texttt{EA-03} & EA & Quais doenças hepáticas são causas de encerramento do tratamento por efeito adverso? & As causas básicas potencialmente relacionadas à hepatotoxicidade de interesse para investigação e registro no IL-TB incluem: Tuberculose (CID A15 a A19), Complicações do HIV (CID B22, B22.7, B23, B23.8), Hepatite aguda medicamentosa (CID K71) e Doença hepática aguda viral (CID B17). \\ \hline
\texttt{EA-04} & EA & O que acontece com o tratamento da ILTB quando a paciente engravida durante o tratamento? & O diagnóstico de gravidez no decorrer do tratamento é uma causa de descontinuidade do tratamento. A paciente deverá iniciar um novo tratamento para ILTB após o parto. Para o novo tratamento, cada caso necessitará ser avaliado individualmente, considerando o número de doses tomadas até a interrupção do tratamento. \\ \hline
\texttt{FE-01} & FE & Qual o tratamento para tuberculose ativa com resistência à rifampicina? & --- \\ \hline
\texttt{FE-02} & FE & Qual a dose de amoxicilina para tratar pneumonia bacteriana? & --- \\ \hline
\texttt{FE-03} & FE & Quais são os critérios diagnósticos para COVID-19? & --- \\ \hline
\texttt{FE-04} & FE & Como fazer o esquema RIPE para tuberculose ativa? & --- \\